

\section{Injective Polynomial Hashing with Fewer Multiplications}
\label{sec:injective}

In this section we use the following recursive family.
\begin{definition}[Injective recurrence]\label{def:injective:recurrence}
Let $u,y,z$ be algebraically independent variables and define
\begin{align}
   P_0 &= z, \\
   P_i &= a_i+(b_i+y)(P_{i-1}+u),
   \qquad i\ge1. \label{eq:injective-recurrence}
\end{align}
\end{definition}
For hashing we use the recurrence with $(u,y,z)$ drawn as three independent
uniformly random field elements, the \emph{three-key family}, or the
\emph{single-key} specialization
\begin{equation}\label{eq:injective-single-key}
   y=x^3,\qquad z=x,\qquad u=x^2,
\end{equation}
that is $P_0=x$ and $P_i=a_i+(b_i+x^3)(P_{i-1}+x^2)$, which keeps a single
resident key word at the price of a weaker bound.
We prove below that distinct parameter vectors
$(a_1,b_1,\dots,a_n,b_n)$ induce distinct polynomials, both before and after
this specialization.  Neither family is bijective onto all monic
polynomials: in the single-key form $P_n$ has degree $3n+2$ but only $2n$
parameters.


\vspace{.5em}

Before proving injectivity, we note the main application: the recurrence gives an almost-universal hash family with half the multiplications of Horner's method.
We use the following standard quantitative definition.
\begin{definition}[Almost-universal hash family]\label{def:injective:universal}
   A family of hash functions $\mathcal H$ is \emph{$\varepsilon$-almost universal} if for any two distinct messages $m\ne m'$ we have
   \[
      \Pr_{h\sim\mathcal H}\!\left[h(m)=h(m')\right] \le \varepsilon.
   \]
\end{definition}

If two distinct messages induce distinct key polynomials whose \emph{difference} has total degree at most $d$, the Schwartz--Zippel lemma gives collision probability at most $d/|\F|$ over uniformly random keys.
In the three-key family $P_n$ is monic of degree $n$ in $y$ with the message-independent top part $(z+u)y^n$ (see the end of the proof of \Cref{lem:injective:coeffmap}), so the difference of two such polynomials has total degree at most $n$, and the family is $n/|\F|$-almost universal on messages of $2n$ field elements: half of Horner's bound at half of Horner's multiplications, with three resident key words.
In the single-key specialization $P_n$ has degree $3n+2$ and its three leading coefficients are again the same for every message, so the difference has degree at most $3n-1$ and the family is $(3n-1)/|\F|$-almost universal with one key word.

Notice a key difference between the bijective polynomials for $k$-wise independent hashing in the previous section, and the injective polynomials here.
In the bijective setting, the polynomial coefficients are the (random) \emph{key} and the input $x$ is the \emph{data} being hashed.
In the injective setting the roles are reversed: the stream $(a_i,b_i)$ is the \emph{data}, while the evaluation point, $x$ is the (random) \emph{key}.

\vspace{.5em}

Before proving the injectivity of our construction, we note the following construction
suggested by Daniel J. Bernstein~\cite{Bernstein:2011:RW}: a simplified version of the
Rabin--Winograd bijective polynomial hashing scheme, which he observed gives an ``authenticator''
\[
   P() = 1, \quad
   P(a_1, \ldots, a_m) = P(a_1, \ldots, a_{k-1})(x^k + a_k) + P(a_{k+1}, \ldots, a_m)
\]
where $k \le m$ is the largest power of $2$ at most $m$.
For example:
$P(a_1)=(x+a_1)+1$,
$P(a_1,a_2) = (x+a_1+1)(x^2+a_2)+1$
and $P(a_1,a_2,a_3) = (x+a_1+1)(x^2+a_2)+x+a_3+1$
and so on.

To see this gives an injective family, first note
that each polynomial (let's write $P_m = P(a_1,\dots,a_m)$ for simplicity)
is monic, since $1$ is monic and product and sums of monic polynomials are monic.
%
Next note that $\deg(P_0)=0$ and $\deg(P_m) = 2k-1$, where $k=k(m)$ is the largest power of 2 less than or equal to $m$ as above.
From the construction, and the fact $m-k<k$, we have:
\[
   \deg(P_m) = \deg(P_{k-1}) + k.
\]
Since $k = 2^j$, we have $k - 1 = 2^j - 1$, so the largest power of 2 $\le k-1$ is $k' = 2^{j-1} = k/2$.
Therefore $\deg(P_{k-1}) = 2(k/2) - 1 = k - 1$,
and $\deg(P_m) = 2k-1$.
%
Finally, the same degree separation gives a decoder.  From the coefficients of
$P_m$ in degrees $k,\ldots,2k-1$ we recover $P(a_1,\ldots,a_{k-1})$, because
only its product with $x^k$ reaches those degrees.  We recursively decode
$a_1,\ldots,a_{k-1}$.  After subtracting the $x^k$ product, the coefficient of
$x^{k-1}$ is $a_k$ plus the known leading coefficient of
$P(a_{k+1},\ldots,a_m)$ when that polynomial has degree $k-1$ (and just
$a_k$ otherwise).  Thus we recover $a_k$, subtract its product with the known
left polynomial, and recursively decode the right polynomial.  This proves
injectivity.

The ``Bernstein Rabin--Winograd'' scheme is simple to describe and prove, but
the recursive nature is a bit inconvenient for fast hashing.
Some authors~\cite{chakrabortyefficient} have proposed non-recursive variants that are more efficient to compute, but they still need to precompute $x^k$ for all powers of 2 up to $n$, which is extra accounting that our method does not require.
Bernstein himself observes that the $\lceil\log_2 m\rceil$ key squarings can be traded for $\lceil\log_2 m\rceil$ cached key powers~\cite[\S5.8]{Bernstein:2011:RW}; the three-key recurrence of \Cref{def:injective:recurrence} needs no key powers at all, and its single-key specialization \eqref{eq:injective-single-key} needs only the two cached powers $x^2,x^3$, for every message length.

\vspace{.5em}



We now prove the injectivity property used by both the one-key and two-key
specializations.

% \begin{lemma}\label{lem:injective:coeffmap}
%    Let $P_i$ be as defined in~\cref{eq:P-def}, and let $c_0, \dots, c_{3n-1}$ denote the coefficients of the monomials $y^k$, $zy^k$, $xy^k$ for $0\le k<n$ in $P_n$.
%    Then the mapping $(a_1, \dots, a_n, b_1, \dots, b_n) \mapsto (c_0, \dots, c_{3n-1})$ is injective.
% \end{lemma}
\begin{lemma}[Injectivity of the recurrence]\label{lem:injective:coeffmap}
The parameter-to-polynomial map in \Cref{def:injective:recurrence} is
injective.  Consequently, the single-variable specialization
\eqref{eq:injective-single-key} is injective as a polynomial in $x$.
\end{lemma}
\begin{proof}
   We first note the following facts about \eqref{eq:injective-recurrence}:
   \begin{enumerate}
      \item $P_n$ has the form $f_1(y) + z f_2(y) + u f_3(y)$,
         where $f_1$ has degree at most $n-1$ and $f_2$ and $f_3$ are degree $n$ polynomials in $y$.
      \item $f_1(y) = \sum_{i=1}^n a_i \prod_{j>i}(y+b_j)$.
      \item $f_2(y) = \prod_{j=1}^n (b_j + y)$,
      \item $f_3(y) = \sum_{i=1}^n \prod_{j\ge i} (b_j+y)$.
   \end{enumerate}
   Here 2. follows if we let $u=z=0$ and expand
   \[P_n=a_n+(b_n+y)P_{n-1}=a_n+(b_n+y)(a_{n-1}+(b_{n-1}+y)P_{n-2}),\]
   etc.
   The other facts follow similarly.

   To show injectivity, we reconstruct the $a$s and $b$s from the coefficients
   of $P_n$.
   It is convenient to start by extracting the $b$s.
   If $n=1$, then $f_2(y)=y+b_1$, so its constant coefficient gives $b_1$
   immediately.  Suppose henceforth that $n\ge2$.
   For this, note that
   \begin{align}
      f_2(y) &= y^n + (\sum_{i=1}^n b_i) y^{n-1} + (\sum_{i<j} b_i b_j) y^{n-2} + O(y^{n-3})
      \quad\text{and}\\
      f_3(y) &= y^{n} + (1+\sum_{i=1}^n b_i) y^{n-1} \\
             &\quad + (1 + \sum_{i=2}^n b_i + \sum_{i<j} b_i b_j) y^{n-2} + O(y^{n-3}),
      \quad\text{so}\\
      f_3(y)-f_2(y) &= y^{n-1} + (1 + \sum_{i=2}^n b_i) y^{n-2} + O(y^{n-3}).
   \end{align}
   (For $n=2$ the constant $1$ in the $y^{n-2}$ coefficient of $f_3$, and hence of $f_3-f_2$, is absent, since it comes from the $i=3$ term of $f_3$; the extraction below is unaffected.)
   We can thus extract $\sum_{i=1}^n b_i$ and $\sum_{i=2}^n b_i$, and from the difference we get $b_1$.
   Using polynomial synthetic division, we may then divide $f_2(y)$ by $(b_1+y)$ to obtain $f_2^{(2)}(y) := \prod_{j=2}^n (b_j + y)$.
   Similarly define $f_3^{(2)}(y) := f_3(y)-f_2(y)=\sum_{i=2}^n \prod_{j\ge i} (b_j+y)$.
   Then proceed by induction to extract the remaining $b$s.

   To extract the $a$s it suffices to notice that $f_1(y) = a_1 y^{n-1} + O(y^{n-2})$ so we may read $a_1$ off the leading coefficient.
   In the process of extracting the $b$s we constructed for each $i$ the polynomial $\prod_{j>i}(y+b_j)$, so we can subtract $a_1\prod_{j>1}(y+b_j)$ and proceed by induction.

   \vspace{.5em}

   The three-variable construction reduces to a \emph{single random key} $x$
   by the specialization \eqref{eq:injective-single-key}.
   The three summands $f_1(x^3)$, $x f_2(x^3)$, and $x^2 f_3(x^3)$ occupy
   distinct exponent classes modulo~$3$, so the specialization preserves injectivity.
   Its degree is $3n+2$.  The coefficients of $x^{3n+2}$, $x^{3n+1}$ and $x^{3n}$ are $1$, $1$ and $0$ for every message: the first two are the leading coefficients of $f_3$ and $f_2$, and no term of exponent $3n$ can arise, since $f_1(x^3)$ has degree at most $3n-3$ while the other two summands only produce exponents $\equiv1,2\pmod 3$.
   Hence two distinct messages give polynomials whose difference has degree at most $3n-1$, and the collision probability is at most $(3n-1)/|\F|$.

   Alternatively, retain $y$ as an independent key and set $z=x$ and $u=x^2$:
   \begin{align}
      P_0 &= x, \\
      P_{i} &= a_{i} + (b_{i}+y)(P_{i-1}+x^2),
   \end{align}
   which has total degree $n+2$ in $(x,y)$; its terms $x^2y^n$ and $xy^n$ have coefficient $1$ for every message, so differences have total degree at most $n+1$ and the collision probability is at most $(n+1)/|\F|$, comparable
   to standard Horner evaluation.
   With all three of $u$, $y$, $z$ drawn as independent uniform keys the bound
   improves further: $f_2$ and $f_3$ are monic of degree $n$, so the
   degree-$(n+1)$ part $zy^n+uy^n$ of $P_n$ is the same for every message, the
   difference of two distinct messages has total degree at most $n$, and the
   collision probability is at most $n/|\F|$ for messages of $n$ pairs, half of
   Horner's $(2n-1)/|\F|$ on the same $2n$ words, with the same three resident
   key words as the single-key form (which keeps $x$, $x^2$ and $x^3$).
\end{proof}

\subsection{Experiments}
\label{sec:injective:experiments}

The recurrence is benchmarked against Horner's rule in
\Cref{sec:experiments:injective}: with the same field arithmetic
($\mathbb F_{2^{64}}$, carryless multiplication) and the same message lengths,
hashing $2N$ words costs $N$ multiplications instead of $2N-1$, and the measured
speedup grows from $1.4\times$ at $2N=8$ to $1.8$--$2.1\times$ at $2N\ge16$ on ARM,
and from $1.1\times$ to $1.9$--$2.3\times$ on x86 (\Cref{tab:injective}).
\Cref{sec:experiments:universal} compares it further with the parallel and
tree-structured universal hashes (Horner-unrolled, BRW, CLNH).

As an implementation data point we also ported the recurrence to
SMHasher3~\cite{smhasher3}, a hash-testing framework, on an Apple M2 Pro
(ARM64), with $L=9$ independent lanes each consuming 16 bytes per step, a
ternary-tree reduction of the lane states and a finalizer that mixes in the
message length.  Over the Mersenne prime field $p=2^{61}-1$ (128-bit product
and two rounds of shift-and-add reduction) it reaches $1.49$~bytes/cycle of
bulk throughput on a 256\,KiB message and $142$~cycles per hash on short
(1--31 byte) keys.  Replacing the field multiplication by the ``MUM'' fold
$\mathrm{lo}_{64}(ab)\oplus\mathrm{hi}_{64}(ab)$ of the $128$-bit product makes
the bulk throughput $9\times$ higher ($14.1$~bytes/cycle), but that variant is
not a field operation: the fold is many-to-one, so the injectivity argument of
this section and the Schwartz--Zippel collision bound do not apply to it, and
message sets with elevated collision rates exist.  We report the number only to
locate the cost of exact field arithmetic; the guarantees of this section are
for the field version.

\paragraph{Polynomial-level comparison.}
\Cref{tab:injective:universal} times the recurrence against the other
$O(1)$-key polynomial methods, and against CLNH as the throughput ceiling of
an $O(N)$-key inner product, on messages of $2N$ words over $\F_{2^{64}}$ with
carryless multiplication, on an Apple M2 Pro and an Intel Xeon Platinum 8375C.
All methods share the same field arithmetic and the same benchmark harness
(\texttt{tools/bench/carryless\_arm.cpp}, \texttt{tools/bench/carryless.cpp}).
Sequentially, the recurrence halves the running time of Horner's rule on both
machines once $2N\ge16$, as its multiplication count predicts.  Fanning out
matters more than the count: the prefix-scan parallelization spends $3N$
multiplications and loses to the tree form of Horner on the Xeon, whereas
dealing the pairs to $L=8$ interleaved lanes keeps the count at $N+L$
(\Cref{sec:experiments:universal}).  With lanes the recurrence is the fastest
$O(1)$-key method on the M2 Pro from $2N=32$ on, by $1.2$--$1.5\times$ over the
next best and within $2\times$ of the $O(N)$-key CLNH ($1.6\times$ at $2N=256$), and on the Xeon from
$2N=128$ on, by $1.2\times$ at $2N=256$; at $32$--$64$ words on the Xeon the
tree form of Horner is $5$--$9\%$ faster, the count advantage being absorbed by
the per-multiplication cost of moving operands between integer and vector
registers in this implementation style (a vector-resident implementation of
the lanes reaches $23.8$\,GB/s on the M2 Pro, \Cref{tab:injective:adversarial}).
BRW is slower than the unrolled, scanned and laned forms on either machine; it
beats only sequential Horner from $2N=32$, the sequential recurrence from
$2N=64$, and, on the M2 Pro, the tree form of Horner up to $2N=128$.

\begin{table}[!htbp]
\centering
\caption{Universal hashing of $2N$ words over $\F_{2^{64}}$, in \textmu s per
$10^5$ hashes (mean over the fastest 100 of 200 repetitions; standard deviations
are below $5\%$ except where marked~$^{\dagger}$).  Horner: sequential, unrolled
over interleaved chains, and the tree evaluation with cached powers
(``parallel'' in \Cref{sec:experiments:universal}).  This paper: the
sequential recurrence ($N$ multiplications), the prefix-scan parallelization
($3N$), and $L=\min(N,8)$ lanes ($N+L$).  CLNH uses $O(N)$ key words; all other
methods use $O(1)$.  Best $O(1)$-key method per row in bold.}
\label{tab:injective:universal}
\resizebox{\columnwidth}{!}{%
\begin{tabular}{c|rrr|rrr|r|r}
 & \multicolumn{3}{c|}{Horner} & \multicolumn{3}{c|}{This paper} & BRW & CLNH \\
 $2N$ & Seq & Unrolled & Tree & Seq & Scan & Lanes & & \\
\hline
\multicolumn{9}{l}{\emph{Apple M2 Pro (ARM, PMULL)}} \\
  8 & 552 & 464 & 2481 & \textbf{393} & 828 & 604 & 1237 & 184 \\
 16 & 2109 & 1045 & 3463 & \textbf{991} & 1544 & 1400 & 2332 & 346 \\
 32 & 7646 & 2727 & 6607 & 4099 & 2735 & \textbf{2107} & 4481 & 1066 \\
 64 & 22976 & 6684 & 12818 & 11672 & 5162 & \textbf{4153} & 8833 & 2115 \\
128 & 64638 & 13710 & 22976 & 27706 & 10288 & \textbf{7850} & 17548 & 4657 \\
256 & 146944 & 34503 & 21157 & 71018 & 22144 & \textbf{14336} & 34879 & 8863 \\
\hline
\multicolumn{9}{l}{\emph{Intel Xeon Platinum 8375C (x86, PCLMULQDQ)}} \\
  8 & 1338 & 1134 & \textbf{1107} & 1192 & 1878 & 1836 & 4302 & 573 \\
 16 & 4560 & \textbf{1999} & 2053 & 2445 & 3168 & 3020 & 6131 & 1153 \\
 32 & 14065 & 7246 & \textbf{4067} & 6002 & 5796 & 4266 & 10350 & 2307 \\
 64 & 36279 & 9140 & \textbf{7443} & 17588 & 10835 & 8135 & 15877 & 4619 \\
128 & 80611 & 20546 & 13725 & 40697 & 20081 & \textbf{13325} & 27669 & 9244 \\
256 & 169985 & 43804 & 26990 & 86905 & 38612 & \textbf{23163} & 47747 & 18492 \\
\end{tabular}
}
\end{table}

\paragraph{Comparison with general-purpose hashes.}
We adopt the standard model for keyed hashing: the secret key material is drawn
uniformly at random and hidden from the attacker, who chooses the input
messages.  The quantity of interest is the worst case, over attacker-chosen
inputs, of the collision probability taken over the random secret.  A universal
hash bounds this by a small $\varepsilon$ for \emph{every} input pair; a
heuristic hash need not.  To compare hashes across message lengths we report
\[
   \mathrm{bits} \;=\; \min_{L}\;\log_2\frac{L}{\varepsilon(L)},
\]
where $L$ is the message length in $64$-bit words and $\varepsilon(L)$ is the
collision probability at that length, floored at $2^{-w}$ for a $w$-bit output.
Dividing by $L$ cancels the Schwartz--Zippel factor, so a degree-$L$ polynomial
hash over a field of size $q$ scores $\log_2 q$ regardless of length, and
taking the minimum keeps the score of a hash with a length-independent flaw
honest.  For the proven hashes $\varepsilon$ is the published bound; for the
heuristic hashes it is the largest collision rate we could \emph{find} by
differential search, measured over $2^{31}$ or more random secrets, so those
entries (marked~$^*$) are upper bounds on the true score: worse inputs may
exist.  We further restrict the main comparison to hashes whose per-key state
is at most eight $64$-bit words, one cache line, following Bernstein's
observation that a server handling thousands of keys cannot keep large per-key
tables hot~\cite{Bernstein:2011:RW}; the budget also makes the throughput
comparison fair, since key size is otherwise a free parameter whose limit is a
key as long as the message.  Shared public constants and per-call derivations
do not count.  Hashes with larger keys are listed below the line.

\Cref{tab:injective:adversarial} gives the result.  Within the budget, the
proven methods score $57$--$64$ bits and the heuristic methods far less:
\texttt{wyhash}, \texttt{rapidhash} and \texttt{XXH3} lose to the all-ones
difference on the two words feeding their first multiply-fold, which collides
with probability about $2^{-27}$ under a random secret, some $2^{37}$ times the
ideal; the $128$-bit \texttt{XXH3-128} fares no better, because one of its
two accumulators sees only a raw word sum that the attacker can hold fixed; \texttt{MUM} has \emph{key-free} collisions, a fixed pair of $8$-byte
messages that collides for every seed; \texttt{komihash} resisted every
differential we tried at $2^{32}$ seeds per surface, which bounds its score
below only by our search effort.  The attacks and their measurements are in
\Cref{app:adversarial}, which also records that the low $64$ bits of an $89$-bit
Mersenne residue are \emph{not} a safe output: a pair differing by $2^{64}-1$
in the last word collides with probability $2^{-25}$, so the Mersenne rows
output the full residue.

\begin{table}[!htbp]
\centering
\caption{Provable security versus throughput on an Apple M2 Pro.
\emph{Key} is the resident per-key state in $64$-bit words.  \emph{Bits} is
$\min_L\log_2(L/\varepsilon(L))$ as defined in the text; entries marked $^*$
are measured on the worst inputs we found and may be lower.  Throughput is
single-threaded, for $16$\,KB and $512$-byte messages.  Rows above the line
keep at most eight words of key.  The two ``this paper'' rows over
$\F_{2^{64}}$ are the three-key recurrence with its state kept in vector
registers (\texttt{tools/bench/adversarial/}); the framework implementation of
\Cref{tab:injective:universal}, which moves the state through general
registers, runs the same one-chain recurrence at $2.2$\,GB/s.  The other
polynomial rows use the framework implementations.  The ChainHash bounds
(\Cref{sec:ph}) are proved for a key of $W+9$ independent uniform field
elements; the implementation timed and tested here expands a $64$-bit seed
with \texttt{splitmix64}, a $2^{64}$-member subfamily that the theorems do
not cover (\Cref{rem:ph:gaps}).}
\label{tab:injective:adversarial}
\begin{tabular}{lrrrr}
\hline
hash & key & bits & GB/s (16\,KB) & GB/s (512\,B) \\
\hline
\multicolumn{5}{l}{\emph{proven, key $\le 8$ words}} \\
This paper, one chain, $\F_{2^{64}}$        & 3 & 64 &  4.1 & 10.1 \\
This paper, 8 lanes, $\F_{2^{64}}$          & 4 & 61 & 23.8 & 17.9 \\
This paper, $\F_{2^{89}-1}$ (89-bit output) & 5 & 88 &  4.4 &  3.9 \\
Horner, $\F_{2^{64}}$                       & 1 & 64 &  1.3 &  2.3 \\
Horner, unrolled                            & 1 & 64 &  5.1 &  7.4 \\
BRW~\cite{Bernstein:2011:RW}                & 1 & 63 &  6.0 &  5.3 \\
Polymur~\cite{polymur}                      & 4 & 57 & 19.7 & 16.2 \\
\multicolumn{5}{l}{\emph{heuristic, key $\le 8$ words}} \\
\texttt{wyhash} v4.3                        & 5 & 28.6$^*$ & 26.5 & 34.8 \\
\texttt{rapidhash} v1                       & 4 & 28.4$^*$ & 27.1 & 32.1 \\
\texttt{XXH3}                               & 1 & 27.3$^*$ & 38.2 & 27.2 \\
\texttt{XXH3-128} (128-bit output)          & 1 & 28.5$^*$ & 33.8 & 24.7 \\
\texttt{MUM} v3                             & 1 &  0$^*$   & 33.0 & 26.8 \\
\texttt{komihash} v5.34                     & 1 & $\ge32^*$ & 24.8 & 24.3 \\
\hline
\multicolumn{5}{l}{\emph{proven, key $> 8$ words}} \\
ChainHash, 1\,KB blocks                     & 137 & 62 & 69.3 & 39.2 \\
ChainHash, 256\,B blocks                    &  41 & 62 & 67.7 & 37.6 \\
ChainHash, 64\,B blocks                     &  17 & 62 & 27.5 & 29.4 \\
UMASH-64~\cite{umash}                       &  38 & 55 & 40.7 & 32.9 \\
UMASH-128 (128-bit output)                  &  38 & 83 & 23.9 & 19.7 \\
CLNH (fixed length)~\cite{lemire2015clhash} & $L$ & 64 & 27.3 & 26.3 \\
Multiply-shift (fixed length)               & $L{+}1$ & 64 & --- & 16.5 \\
\hline
\end{tabular}
\end{table}

\paragraph{A large-key hash built from the recurrence.}
The rows labelled ChainHash compose the recurrence with the reduction-free
inner level of the NH-based hybrids, in the same way that UMAC, VMAC and CLHASH
compose a polynomial with NH~\cite{black1999umac,lemire2015clhash}.  The name
reflects the structure: each step of the recurrence is an NH-style product into
which the running value is mixed, and that mixing is what lets one key be reused
across the whole message instead of one key word per message word.  The message
is split into blocks of $B$ bytes and each block into $S$ sub-blocks ($S=1$ for
$B\le256$, $S=2$ for $1$\,KB blocks); each sub-block is hashed by carryless NH
(one \texttt{PMULL} per $16$ bytes into two $128$-bit accumulators, the words
of every $32$-byte group paired as $(w_{4j},w_{4j+2})$ and
$(w_{4j+1},w_{4j+3})$, which is what two vector loads deliver without a
shuffle) and its
unreduced sum is consumed as one pair $(a_i,b_i)$ of the three-key recurrence
$P_0=z$, $P_i=a_i+(b_i+y)(P_{i-1}+u)$, one field multiplication per
sub-block; the byte length is XORed into both halves of the last pair; one further key
word is added to the result modulo $2^{64}$ (the \emph{twist}: an integer
addition, carries and all); and the sum is passed through a random monic
polynomial of degree~$5$ evaluated with three multiplications,
$\lfloor5/2\rfloor+1$ as in \Cref{thm:main}, by a characteristic-$2$ circuit
whose coefficient map is a bijection of $\F^5$ (\Cref{sec:ph} gives the
explicit decoder).  With its five parameters uniform, the finalizer is a
$5$-wise independent family on distinct field inputs; for the full hash
\Cref{thm:ph:kwise} gives the qualified form: conditioned on the level-2
values of up to five messages being distinct, which fails with probability
at most $\binom{t}{2}(p+1)/2^{64}$, the outputs are independent and uniform,
so that two messages collide with probability $r+(1-r)/2^{64}$, $r$ being the
level-1/2 collision probability, rather than $1/2^{64}$.  Five-wise
independence is what linear probing provably
needs~\cite{pagh2009linear,patracscu2016kindependence}.  Both statements are
for the ideal key of $W+9$ uniform words; the implementation derives its key
from a $64$-bit seed (\Cref{rem:ph:gaps}).
The twist answers a question that $k$-wise independence never asks: over
$\F_{2^{64}}$ the $\F_2$-degree of $v\mapsto v^e$ is the number of ones in
the binary expansion of $e$, so every polynomial of degree at most $6$ is
at most quadratic in the bits of its input, with affine discrete derivatives, and the
fixed-seed keysets of SMHasher3~\cite{smhasher3} (Zeroes, Sparse,
Permutation, TwoBytes, Bitflip) are sensitive to exactly that structure: the
untwisted degree-$5$ finalizer fails $22$ of the $200$ tests, and an earlier
version of the hash used degree $7$, the first cubic option, at the price of
a fourth multiplication.  Any fixed bijection of the finalizer input
preserves the $2^{-64}$ bound and the $5$-wise independence exactly, and the
carry chain of an integer addition is not $\F_2$-affine
(\Cref{sec:ph:twist}); with the twist the degree-$5$ hash passes $200/200$,
degree $3$ with the twist still fails $17$, and whether the twist suffices
for a given test suite is an empirical question, not a theorem.  $S=2$
gives a one-block message a second message-dependent multiplication for the
same kind of reason.  The collision probability is at most $(p+2)/2^{64}$ for two
messages of at most $n$ blocks, $p=Sn$ sub-blocks: $2^{-64}$ for the block
level (the full-width XOR-universality of carryless NH, extended to last
blocks of different pair counts by the length term), $p/2^{64}$ for the
recurrence (the three-key bound above, with $p$ counting sub-blocks rather
than word pairs) and $2^{-64}$ for the finalizer (\Cref{sec:ph}).  No
heuristic mixing step is used anywhere, and both variants pass all $200$
tests of the full SMHasher3 suite, on an Apple M2 Pro and on an Intel Xeon
8375C, including a new seeded-differential test
(full-output collisions of structured pairs, complemented leading or adjacent
interior words, over $2^{24}$ random seeds, $2^{30}$ in an extended tier) that
\texttt{MUM} fails outright and \texttt{wyhash}, \texttt{rapidhash},
\texttt{XXH3} and \texttt{XXH3-128} fail in the extended tier, while \texttt{komihash},
Polymur, SipHash and the polynomial hashes pass (\Cref{app:adversarial}).
With $1$\,KB blocks and a $137$-word key it runs at $69$\,GB/s on the M2 Pro,
above every heuristic hash in the table and above UMASH, with a proof---for
the ideal key; the seeded interface that was timed and tested is a subfamily
outside the theorems (\Cref{rem:ph:gaps}); the
block size is exactly the key-budget knob, and at $64$-byte blocks, closest to
the one-cache-line budget ($17$ key words), the hybrid matches or beats the
eight-lane recurrence on long messages ($27.5$ against $23.8$\,GB/s at
$16$\,KB).  The implementation and a
bit-serial reference are in \texttt{tools/bench/chainhash/} of the repository.

\paragraph{Key expansion.}
One may ask whether the $O(B)$ key of the block level can be derived from a few
seed words, say by a $k$-wise independent polynomial.  The collision condition
of NH is affine in the key words, because the products of two key words cancel
in the difference of two block sums, so with polynomially expanded keys it is a
linear map from the seed words to the difference.  For a \emph{reduced} NH,
whose block sums are taken in $\F_{2^{64}}$, this is fatal: a difference
supported on $k+1$ positions chosen in the kernel of the $k\times(k+1)$
Vandermonde system is independent of the key, and for $k=6$ we verified that
such a pair of $38$-word blocks has the same difference under $200$ random
seeds, a key-free collision.  For the \emph{unreduced} carryless NH used here
the products are taken in $\F_2[X]$ while the expansion lives in
$\F_{2^{64}}$, and the same crafted difference only lowers the $\F_2$-rank of
the seed-to-difference map from $127$ to $63$, leaving a collision probability
of about $2^{-63}$.  We found no better attack, but also no proof; the
security of unreduced NH under $k$-wise independent keys is an open problem
(\Cref{sec:open-problems}).  UMAC's answer is to expand the key with AES,
trading the information-theoretic bound for a computational one.

